% ===== airbnb-frames.tex: the Airbnb case in three frames, asked for by Carlos on 2026-11-06 =====
% The frames of the case and their talk track. Since 2026-11-09 two decks input this file, each as one section:
% block1-nearest-neighbors.tex, the deck of its Friday block, and weekend2-cases.tex, the weekend deck.
% His dictated text: (1) the first stage of the search evaluates tens of thousands of homes and passes a shortlist
% to the ranking model; it used to be a simple linear model that scored every home once a day; (2) one network
% turns the search into a list of numbers, a second turns each home into a list of numbers in a daily batch job, and
% the homes nearest to the search join the shortlist; (3) each version A/B tested, +0.31% conversion together, "a
% large portion of the team's annual goal", on par with the largest gains of two years (blog), +2.3% bookings from
% promotional emails, about 16% less compute with the clustered index. Diagrams as for the other case decks.
% Simplified the same day after his "This is too complicated for me ... Make them 25% simpler": shorter text and
% labels, fewer points in frame 2, one quote in frame 3 (the blog's comparison is in the talk track). README.md,
% section "One deck per case", has the counts.
% Since 2026-11-09 frame 3 names its measure by what it counts, "visits that end in a booking", at Carlos's request
% ("Please relabel as you suggest"); the paper says "booking conversion" and does not define it. Since 2026-11-10
% the paper's words on that gain are introduced with his sentence, "Airbnb calls this gain"; they belong on this
% frame and nowhere in another case (his words: "They need to go in the Airbnb part").
% No number is typed here: numbers.tex is written by tools/cases_numbers.py (block Airbnb), which keeps the sentence
% of each source. The points of frame 2 are a drawing, not data. A "Say" comment above each frame is the talk
% track. Build and check from this folder:
%   python3 tools/cases_numbers.py
%   ../style/tools/build.sh block1-nearest-neighbors.tex
%   python3 tools/check_numbers.py block1-nearest-neighbors.pdf tools/expected-block1-nearest-neighbors.txt
\newcommand{\absource}{Abdool and others (Airbnb), arXiv preprint 2026; Airbnb's blog 2025.}

% ===== Frame 1. The first stage of Airbnb's search =====
% Say: Airbnb, the online marketplace for places to stay. A popular location can have around a hundred
% thousand available homes. The first stage of its search evaluates tens of thousands of them and passes a
% shortlist to the ranking model, which orders it. [click] That first stage used to be only a simple linear model
% that scored every home once a day. One score per home, whatever you searched: among the homes that fit both
% searches, the couple and the family of six got the same order.
\begin{frame}{Airbnb: the first stage of search}
  \begin{tikzpicture}[x=1mm, y=1mm]
    \path[use as bounding box] (0,1) rectangle (130,-54);
    \node[body, anchor=north west, text width=64mm] at (0,0) {The first stage evaluates tens of thousands of homes.};
    \node[body, anchor=north west, text width=64mm, visible on=<2->] at (0,-11)
      {It used to be only a simple linear model that scored every home once a day, whatever the search.};
    \node[stage, minimum width=59mm, text width=55mm, align=center] (all) at (99.5,-5.6)
      {In a popular place, about\\ a hundred thousand homes};
    \node[stage, minimum width=59mm] (short) at (99.5,-21) {A shortlist};
    \node[stage, minimum width=59mm] (ranked) at (99.5,-35) {Ranked results};
    \draw[flow] (all) to node[elab, inner sep=0.4pt, anchor=west, xshift=1.2mm] {first stage} (short);
    \draw[flow] (short) to node[elab, inner sep=0.4pt, anchor=west, xshift=1.2mm] {ranking model} (ranked);
    \node[src, anchor=north west, text width=130mm] at (0,-50.6) {\absource};
  \end{tikzpicture}
\end{frame}

% ===== Frame 2. Two networks and the nearest homes =====
% Say: What Airbnb added to the first stage. One network turns the search, its location, the number of guests
% and the size of the area, into a list of numbers, at the moment of the search. A second network turns each
% home, its amenities, location and capacity, its past views, wishlists and reviews, into a list of numbers of the
% same length, in a batch job once a day for every home. [click] Picture every home as a point and the search as
% a point: the homes nearest to the search join the shortlist. In the first version "nearest" meant the highest
% inner product, as in our video store; the latest one uses the distance.
\begin{frame}{Two networks and the nearest homes}
  \begin{tikzpicture}[x=1mm, y=1mm]
    \path[use as bounding box] (0,1) rectangle (130,-54);
    \node[anchor=north, inner sep=0, font=\small, align=center] at (19,-2) {Location,\\ number of guests};
    \node[anchor=north, inner sep=0, font=\small, align=center] at (63,-2) {Amenities,\\ location, reviews};
    \node[net] (snet) at (19,-19) {Search network};
    \node[net] (hnet) at (63,-19) {Home network};
    \draw[flow] (19,-10.4) to (snet.north);
    \draw[flow] (63,-10.4) to (hnet.north);
    \draw[flow] (snet.south) to (19,-27.3);
    \draw[flow] (hnet.south) to (63,-27.3);
    \hatchcells{1-}{12.1}{-29.5}{4}
    \hatchcells{1-}{56.1}{-29.5}{4}
    \node[lbl, anchor=north] at (19,-32.4) {at each search};
    \node[lbl, anchor=north] at (63,-32.4) {once a day};
    % the homes as points and the search (click 2); a drawing, the coordinates are not data
    \node[lbl, anchor=north west, text width=40mm, visible on=<2->] at (90,-2) {The nearest homes join the shortlist};
    \foreach \hx/\hy in {94/-14, 118/-15, 126/-23, 92/-27, 97/-40, 116/-42, 125/-36, 101/-17, 120/-30}
      {\draw[draw=nib, line width=0.6pt, fill=paper, visible on=<2->] (\hx,\hy) circle (1.05);}
    \foreach \hx/\hy in {104.6/-25.2, 112.2/-26.6, 108/-32.4}
      {\fill[hatch fine, visible on=<2->] (\hx,\hy) circle (1.05);
       \draw[draw=brand, line width=0.7pt, visible on=<2->] (\hx,\hy) circle (1.05);}
    \draw[draw=brand, line width=0.8pt, visible on=<2->] (108.4,-28.3) circle (7.2);
    \fill[brand, visible on=<2->] (108.4,-28.3) circle (1.25);
    \node[elab, anchor=south, inner sep=0.6pt, visible on=<2->] at (108.4,-18.6) {search};
    \node[lbl, anchor=west, visible on=<2->] at (99,-40) {a home};
    \node[src, anchor=north west, text width=130mm] at (0,-50.6) {\absource};
  \end{tikzpicture}
\end{frame}

% ===== Frame 3. The value =====
% Say: Airbnb built it in three versions and A/B tested each one: first the two networks with the inner product;
% then the same with a clustered index, where the homes are grouped every day and only the groups nearest to the
% search are searched; then a larger network that uses the distance. [click] Together the three raised the share
% of visits that end in a booking by 0.31 percent, a relative change. That sounds small. Airbnb calls this gain
% "a large portion of the team's annual goal", and its blog of April 2025 put the bookings lift of the new system
% on par with some of its largest machine learning improvements to search ranking in the past two years.
% [click] The same system now also
% picks the homes for promotional emails, where bookings rose by 2.3 percent, and the clustered index used about
% 16 percent less compute than the old system. The source is a preprint of January 2026 and Airbnb's blog of
% April 2025.
\begin{frame}{The value: three A/B tests}
  \begin{tikzpicture}[x=1mm, y=1mm]
    \path[use as bounding box] (0,1) rectangle (130,-54);
    \foreach \vx/\vname/\vtext in {20/v1/{Two networks}, 65/v2/{A clustered index}, 110/v3/{A larger network}}
      {\node[chip, font=\small, minimum width=38mm, minimum height=7.6mm, inner sep=1.5pt] (\vname) at (\vx,-4) {\vtext};}
    \draw[flow] (v1) to (v2);
    \draw[flow] (v2) to (v3);
    \draw[draw=ruledark, line width=0.6pt, line join=miter, visible on=<2->] (1,-9.4) to (1,-10.8) to (129,-10.8) to (129,-9.4);
    \draw[draw=ruledark, line width=0.6pt, visible on=<2->] (22,-10.8) to (22,-12.2);
    \node[big, anchor=north west, visible on=<2->] at (0,-14.6) {+\abConv\%};
    \node[measure, anchor=north west, text width=60mm, visible on=<2->] at (0,-22.4)
      {visits that end in a booking, all three, relative change};
    \node[body, anchor=north west, text width=60mm, visible on=<2->] at (70,-14.6)
      {Airbnb calls this gain \textquotedblleft a large portion of the team's annual goal\textquotedblright.};
    \draw[draw=ruling, line width=0.4pt, visible on=<3->] (0,-33) to (130,-33);
    \node[big, anchor=north west, visible on=<3->] at (0,-36.6) {+\abEmail\%};
    \node[measure, anchor=north west, text width=36mm, visible on=<3->] at (24,-37.2) {bookings from promotional emails};
    \node[big, anchor=north west, visible on=<3->] at (70,-36.6) {\textminus\abCompute\%};
    \node[measure, anchor=north west, text width=38mm, visible on=<3->] at (91,-37.2) {compute with the clustered index};
    \node[src, anchor=north west, text width=130mm] at (0,-50.6) {\absource};
  \end{tikzpicture}
\end{frame}

